\documentclass{article}
\usepackage{colm2024_conference}
\usepackage{microtype}
\usepackage{hyperref}
\usepackage{url}
\usepackage{booktabs}
\usepackage{graphicx}
\usepackage{amsmath}
\usepackage{array}

\definecolor{darkblue}{rgb}{0,0,0.5}
\hypersetup{colorlinks=true,citecolor=darkblue,linkcolor=darkblue,urlcolor=darkblue}

\title{GPT-2 Implementation and Emotion-Conditioned\\
Bilingual Poetry Generation}

\author{Jihong Park \quad Jeewu Lee \quad Hyeongwon Jo \\
Sogang University}

\colmfinalcopy

\begin{document}
\maketitle
\fancyhead{}
\lhead{CSEG321 Final Project Report}

\begin{abstract}
We implement GPT-2 and evaluate it on sentiment classification, paraphrase detection, and Shakespearean sonnet continuation. Full-model fine-tuning reaches development accuracies of 0.514 on SST-5, 0.971 on CFIMDB, and 0.896 on Quora Question Pairs, outperforming frozen-backbone linear probes in sentiment classification. We then test whether explicit emotion information improves English and Korean poem continuation. Our final oracle experiment supplies the complete gold line-level emotion trajectory, including future lines, during training and inference. Across three English sentiment measurements and two Korean label resolutions, conditioned models do not consistently surpass the unconditioned baseline or shuffled controls. An average trajectory sometimes improves the emotional-turn location without improving textual similarity. Thus, in our setting, text-prefix trajectories do not reliably ground emotion control in GPT-2 generation.
\end{abstract}

\section{Introduction}

GPT-2 is a decoder-only Transformer language model whose causal objective supports both text generation and task-specific adaptation~\citep{radford2019language,vaswani2017attention}. In the first part of this project, we implement its major architectural components and apply the model to three downstream tasks: sentiment classification, paraphrase detection, and sonnet generation. This provides a controlled basis for examining both discriminative adaptation and open-ended generation.

The second part connects sentiment analysis and generation. Initial experiments with a global label, sequential emotion fine-tuning, and prefix-only trajectories did not improve reference-based quality. To test whether the condition was too sparse or coarse, we conduct an upper-bound experiment that exposes the model to the complete gold emotion trajectory, including future continuation lines unavailable in deployment.

We implement and validate GPT-2's attention, Transformer block, embeddings, pretrained-weight mapping, and AdamW; establish three downstream baselines; and compare gold, average, and shuffled bilingual emotion trajectories. Even oracle trajectories fail to improve continuation consistently, implicating text-prefix conditioning rather than merely insufficient labels.

\section{GPT-2 implementation and base tasks}

\subsection{Model implementation}

Our implementation follows GPT-2 small (approximately 124M parameters). Each Transformer layer uses pre-layer normalization, masked multi-head self-attention, a position-wise feed-forward network, and residual connections. A triangular causal mask prevents attention to future tokens, while a padding mask excludes padded positions. Learned token and positional embeddings are summed before the Transformer stack, and a final layer normalization produces contextual representations. The language-model head shares weights with the token embedding matrix. We additionally implement decoupled weight decay in AdamW and map pretrained GPT-2 weights into our module naming and tensor layout. The supplied numerical sanity checks pass after implementation.

For Windows and recent PyTorch, we specify UTF-8 input, load checkpoints with the compatible \texttt{weights\_only} setting, and use UTF-8 console output. The starter paraphrase head was incompatible with its ``yes''/``no'' token-ID targets, so we use the shared vocabulary projection and compute cross-entropy over the GPT-2 vocabulary.

\subsection{Tasks and training protocol}

SST-5 is a five-class sentence-level sentiment task~\citep{socher2013recursive}; CFIMDB is a binary counterfactual movie-review task~\citep{kaushik2020learning}. We classify the final non-padding representation with a linear head and compare full-model fine-tuning against training only that head. Paraphrase detection uses Quora Question Pairs and predicts a verbalized binary answer through the language-model head. The final base task is Shakespearean sonnet continuation: given the first three lines, the model generates the remaining eleven lines. The dataset contains 131 training sonnets and 12 sonnets each for development and held-out evaluation.

All runs use our AdamW implementation for 10 epochs with seed 11711. Full-model runs use a learning rate of $10^{-5}$; linear probes use $10^{-3}$. Batch sizes are 64 for SST, 8 for CFIMDB and sonnet generation, and 32 for Quora. We retain the checkpoint with the highest development accuracy for classification. Sonnet generation uses nucleus sampling; the original task has no official generation metric, so we report chrF against the reference continuation~\citep{popovic2015chrf}.

\begin{table}[h!]
\caption{Development results on the three base downstream tasks. Accuracy is the checkpoint-selection metric for classification.}
\label{tab:base-results}
\centering
\small
\begin{tabular}{llrr}
\toprule
Task & Fine-tuning & Accuracy & F1 \\
\midrule
SST-5 & full model & \textbf{0.514} & \textbf{0.496} \\
SST-5 & last linear & 0.445 & 0.400 \\
CFIMDB & full model & \textbf{0.971} & \textbf{0.971} \\
CFIMDB & last linear & 0.857 & 0.856 \\
Quora pairs & full model & \textbf{0.896} & 0.889 \\
Sonnet continuation & full model & \multicolumn{2}{c}{chrF: 27.7--27.8} \\
\bottomrule
\end{tabular}
\end{table}

Full fine-tuning improves accuracy by 0.069 on SST-5 and 0.114 on CFIMDB (Table~\ref{tab:base-results}) and is more stable across epochs than linear probing. Quora reaches 0.896 accuracy and 0.889 F1 after correcting the output-target mismatch. Sonnet training loss decreases from 5.047 to 4.065; outputs reproduce archaic phrasing but have limited long-range coherence.

\section{Emotion-conditioned poetry generation}

\subsection{Research question and progression}

Our main question is whether emotion supervision improves GPT-2 poem continuation. Earlier experiments tested (i) a predicted global emotion label prepended to the poem, (ii) a random-label control, (iii) sequential fine-tuning on emotion classification followed by poetry, and (iv) trajectories derived only from observed prefix lines. Their chrF scores were tied with or below the baseline. We therefore design a final upper-bound test: if a model receives the correct emotion of every line, can it produce a continuation closer to the reference and follow the intended emotional development?

This setting intentionally leaks information from future lines. It is not a deployable method; it estimates the maximum benefit available to the conditioning mechanism. If even this oracle cannot help, improving a prefix emotion classifier alone is unlikely to solve the problem.

\subsection{English and Korean data}

For English, we use the Shakespeare continuation split and represent each 14-line sonnet by three alternative sentiment trajectories from prior computational analysis of the Shakespearean \emph{volta}, or emotional turn~\citep{lee2025volta}. The measurements are denoted $m$ (discrete human labels), $v$ (lexicon-based continuous scores), and $r$ (neural-model continuous scores). Their moderate pairwise correlations make them useful for testing whether conclusions depend on one sentiment instrument. A duplicate was found between development sonnet 138 and the nominal held-out set, so final evaluation excludes it ($n=11$). Conditioned models use 130 full 14-line training sonnets and the duplicate-filtered development set; the unconditioned baseline retains the original 131-poem training and 12-poem development splits for checkpoint selection.

For Korean, we use KoGPT2 and KPoEM modern Korean poems. The generation input is the first 30\% of lines rather than exactly three lines because poem length varies greatly. Training uses 368 poems, and final evaluation uses 11 poems by Kim Sowol, the subset that had shown the most favorable signal in a preceding experiment. We test both the original 44-category line labels and a manually collapsed six-class representation. Korean conditions use emotion names rather than numeric IDs so the pretrained tokenizer can exploit their lexical meaning.

\subsection{Conditions and evaluation}

The unconditioned \textsc{base} receives only poem text. \textsc{Oracle} receives the poem's gold full trajectory. \textsc{Average} receives the mean trajectory over training poems, testing whether a genre-level structural prior helps without poem-specific information. \textsc{Shuffled} receives another poem's real trajectory under a deterministic cyclic permutation. This is a stronger control than independent random values because it preserves the marginal distribution and temporal smoothness of real trajectories while destroying poem--trajectory correspondence.

All final models use full-model fine-tuning for 10 epochs, batch size 8, learning rate $10^{-5}$, and best development loss. Maximum sequence length is 384 for conditioned English and 768 for conditioned Korean. At inference, we generate three deterministic-seed samples per poem with temperature 0.8, top-$p$ 0.9, and at most 200 new tokens, then average sample scores.

Our primary text metrics are chrF and BERTScore F1~\citep{popovic2015chrf,zhang2020bertscore}. We additionally measure trajectory correlation, Distinct-2, \emph{volta distance}, and, for Korean, line-level emotion agreement. The same emotion measurement model is applied to generations and references within each language.

\paragraph{Metric definitions.}
All text-similarity scores compare the generated continuation $y$ with the reference continuation $y^*$. chrF computes an $F_{\beta}$ score from character $n$-gram precision $P_c$ and recall $R_c$ (orders 1--6, $\beta=2$ in SacreBLEU):
\begin{equation}
\operatorname{chrF}_{\beta}=100\frac{(1+\beta^2)P_cR_c}{\beta^2P_c+R_c}.
\end{equation}
BERTScore aligns contextual token embeddings between $y$ and $y^*$ by maximum cosine similarity and reports the harmonic mean of embedding-based precision and recall. Higher chrF and BERTScore indicate greater surface and semantic similarity, respectively.

Let $s=(s_1,\ldots,s_n)$ and $s^*=(s^*_1,\ldots,s^*_m)$ be measured line-level emotion trajectories for the generation and reference. Trajectory correlation is Pearson's $r$ over the first $L=\min(n,m)$ continuation lines; higher values indicate more similar emotional development. The volta position and distance are
\begin{equation}
v(s)=\arg\max_{1\leq i<n}|s_{i+1}-s_i|,\qquad
d_v=|v(s)-v(s^*)|,
\end{equation}
computed on the complete poem (observed prefix plus continuation); lower $d_v$ is better. Distinct-2 is the number of unique whitespace-token bigrams divided by the total number of generated bigrams, so higher values indicate less repetition. Korean line-emotion agreement is $L^{-1}\sum_{i=1}^{L}\mathbf{1}[\ell_i=\ell_i^*]$ over aligned continuation lines, where $\ell_i$ is the predicted emotion label. Per-poem scores are first averaged across the three generated samples and then across poems.

\section{Results}

\subsection{English trajectories}

\begin{table}[h!]
\caption{English sonnet results ($n=11$, three samples per poem). VDist is volta distance (lower is better); Corr is trajectory correlation.}
\label{tab:english}
\centering
\scriptsize
\setlength{\tabcolsep}{4.5pt}
\begin{tabular}{lrrrrr}
\toprule
Condition & chrF & BERT-F1 & VDist$\downarrow$ & Corr & Dist-2 \\
\midrule
Base & \textbf{27.69} & \textbf{0.833} & 4.58 & -0.020 & \textbf{0.840} \\
Oracle $m$ & 25.29 & 0.827 & 4.76 & 0.019 & 0.803 \\
Average $m$ & 24.88 & 0.826 & 3.52 & 0.012 & 0.775 \\
Shuffled $m$ & 27.08 & 0.827 & 3.97 & 0.024 & 0.805 \\
Oracle $v$ & 25.39 & 0.827 & 4.52 & -0.146 & 0.797 \\
Average $v$ & 25.86 & 0.828 & \textbf{2.94} & -0.028 & 0.792 \\
Shuffled $v$ & 26.14 & 0.826 & 5.12 & 0.057 & 0.785 \\
Oracle $r$ & 26.50 & 0.830 & 4.27 & 0.058 & 0.808 \\
Average $r$ & 26.48 & 0.827 & 4.58 & -0.090 & 0.786 \\
Shuffled $r$ & 23.87 & 0.824 & 3.70 & \textbf{0.108} & 0.767 \\
\bottomrule
\end{tabular}
\end{table}

The unconditioned model ranks first in both chrF and BERTScore (Table~\ref{tab:english}). Relative to base, oracle $m$ and $v$ reduce mean per-poem chrF by 2.40 and 2.30 and are worse on 10 of 11 poems; oracle $r$ reduces it by 1.19 and is worse on 8 of 11. Oracle conditions also reduce Distinct-2, indicating that the added prefix tends to increase repetition rather than useful control.

Correct trajectories do not consistently beat content-destroying shuffled controls. Oracle $m$ is 1.79 chrF below shuffled $m$; oracle $v$ is 0.76 below shuffled $v$; oracle $r$ is 2.63 above shuffled $r$, but this is driven by shuffled $r$ being the worst condition overall, and oracle $r$ still trails base. The direction changes across sentiment instruments, so the evidence does not support reliable use of poem-specific trajectory content.

The only repeated positive effect is structural. Average $v$ reduces volta distance from 4.58 to 2.94, and average $m$ reduces it to 3.52, despite both lowering textual similarity. A single mean trajectory appears learnable as a genre prior for where a sonnet tends to turn, whereas poem-specific values do not ground line-level generation.

\subsection{Korean trajectories}

\begin{table}[h!]
\caption{Korean results on 11 Kim Sowol poems with 30\% line input. Agree is line-level emotion agreement; Corr has nine valid poems.}
\label{tab:korean}
\centering
\small
\begin{tabular}{lrrrrr}
\toprule
Condition & chrF & BERT-F1 & VDist$\downarrow$ & Corr & Agree \\
\midrule
Base & 6.28 & 0.612 & 2.30 & \textbf{0.156} & \textbf{0.326} \\
Oracle 44-class & \textbf{6.50} & 0.614 & \textbf{1.73} & -0.025 & 0.248 \\
Oracle 6-class & 6.42 & \textbf{0.618} & 2.27 & 0.041 & 0.287 \\
\bottomrule
\end{tabular}
\end{table}

The Korean text metrics show small numerical gains: the 44-class condition gains 0.22 chrF, and the six-class condition gains 0.006 BERTScore. Given only 11 poems and one training seed, these magnitudes are not sufficient evidence of a robust quality improvement. More importantly, both oracle models follow the target emotion trajectory less closely than base. Correlation falls from 0.156 to $-0.025$ and 0.041, while line-level agreement falls from 0.326 to 0.248 and 0.287. Collapsing 44 labels to six mitigates some degradation but does not recover baseline emotion tracking. Thus, finer label resolution is not the primary bottleneck.

\section{Discussion}

The results suggest a grounding failure. Loss is assigned to both the regular trajectory block and poem text, so a model can learn the block without grounding its $i$th value in the $i$th generated line. Conditioned development losses are therefore not comparable with base loss. In Korean, direct emotion-following metrics worsen even with the gold sequence.

Data scale compounds the problem: 130 conditioned English and 368 Korean poems must teach a 124--125M parameter model a new position-to-line alignment. The average trajectory, by contrast, repeats one pattern across all examples, explaining its learnable structural effect.

Reference metrics may penalize valid alternative wording, but correlation and line-level agreement also fail to improve. Future work should mask condition-block loss, interleave labels with target lines, add an emotion-consistency objective, or guide decoding with a classifier. Larger backbones and multiple seeds are also needed.

\section{Limitations}

Final evaluation uses 11 poems per language and one training seed. The English baseline has one additional training poem and uses the original 12-poem development split, unlike duplicate-filtered conditioned runs. Oracle trajectories deliberately leak future information. The Korean measurer has a comment-to-poetry domain mismatch, valence mapping discards nuance, and 200-token decoding can truncate long targets. Automatic metrics also omit meter, rhyme, imagery, and human preference.

\section{Conclusion}

We implemented GPT-2 and adapted it to three downstream tasks; full fine-tuning outperformed sentiment linear probes. Across global labels, sequential transfer, partial trajectories, and complete oracle trajectories, emotion conditioning provides no reliable improvement. Gold trajectories are not systematically better than shuffled ones, and Korean oracle conditions reduce direct emotion tracking. Text-prefix conditioning was therefore ineffective in our setting; future work should change how conditions connect to generation rather than only adding label detail.

\bibliography{report}
\bibliographystyle{colm2024_conference}

\clearpage
\appendix
\section{Experimental details}

\paragraph{Optimization and decoding.}
Table~\ref{tab:hyperparameters} records the effective settings used for the reported runs. All random initialization, data loading, and generation start from seed 11711. For sample $k$ of poem $i$, decoding uses seed $11711+1000k+i$, giving matched sampling randomness across setups.

\begin{table}[h!]
\caption{Hyperparameters used in the base tasks and final trajectory experiments.}
\label{tab:hyperparameters}
\centering
\scriptsize
\setlength{\tabcolsep}{4pt}
\begin{tabular}{lll}
\toprule
Component & Setting & Value \\
\midrule
Base classification & epochs / optimizer & 10 / AdamW \\
Full-model classifiers & learning rate & $10^{-5}$ \\
Linear probes & learning rate & $10^{-3}$ \\
SST / CFIMDB / Quora & batch size & 64 / 8 / 32 \\
Trajectory SFT & epochs / batch / learning rate & 10 / 8 / $10^{-5}$ \\
Checkpointing & selection & best development accuracy or LM loss \\
Conditioned length & English / Korean & 384 / 768 tokens \\
Decoding & temperature / top-$p$ & 0.8 / 0.9 \\
Generation & samples / maximum new tokens & 3 / 200 \\
Korean decoding & no-repeat / repetition penalty & 3-gram / 1.3 \\
\bottomrule
\end{tabular}
\end{table}

\paragraph{Prefix representation.}
English conditions prepend a textual block to the observed poem:
\begin{quote}\small\ttfamily
\lbrack sentiment trajectory\rbrack\quad
1: -0.169; 2: -0.123; $\ldots$; 14: +0.238\quad
\lbrack/ sentiment trajectory\rbrack\\
FIRST-THREE-POEM-LINES
\end{quote}
The Korean format is analogous but uses Korean emotion names and the first 30\% of poem lines. During SFT, cross-entropy is applied to both the trajectory block and the full poem, with padding labels set to $-100$. At inference, the same block precedes only the observed prefix. Oracle blocks contain the gold trajectory of all lines; average blocks repeat the training-set mean trajectory; shuffled blocks cyclically assign another poem's real trajectory.

\paragraph{Experiment progression.}
Table~\ref{tab:progression} summarizes why the final upper-bound experiment was needed. Values are continuation chrF; differences below approximately 0.1 in the Korean intermediate runs are treated as practically tied given the small evaluation sets.

\begin{table}[h!]
\caption{Progression from coarse emotion supervision to full-trajectory oracle conditioning.}
\label{tab:progression}
\centering
\scriptsize
\begin{tabular}{llll}
\toprule
Stage & Language & Compared setups & chrF \\
\midrule
Global / transfer & EN & base / predicted / random / sequential & 27.81 / 26.57 / 26.41 / 23.67 \\
Global / transfer & KO & base / predicted / random / sequential & 7.16 / 7.17 / 7.27 / 7.26 \\
Early trajectory & EN & base / oracle-$m$ / average-$m$ & 27.81 / 24.96 / 24.75 \\
Early trajectory & KO & base / oracle / average / random & 7.16 / 7.13 / 7.09 / 7.08 \\
Full trajectory & EN & base / best oracle & 27.69 / 26.50 \\
Full trajectory & KO & base / 44-class / 6-class & 6.28 / 6.50 / 6.42 \\
\bottomrule
\end{tabular}
\end{table}

\paragraph{Measurement details.}
English line sentiment uses the SST-2 DistilBERT model (\nolinkurl{distilbert-base-uncased-finetuned-sst-2-english}) and score $P(\text{positive})-P(\text{negative})$. Korean lines use our six-class KOTE classifier with joy/love $=+1$, surprise $=0$, and sadness/anger/fear $=-1$. The same measurer scores generations and references. BERTScore selects its language model via \texttt{lang=en/ko}; exact reruns should pin that backbone. English trajectory metrics have 11 valid poems. Korean text, volta, and agreement metrics use 11; correlation is defined for nine.

\end{document}
